\documentclass[11pt]{article}
\usepackage[margin=1in]{geometry}
\usepackage{amsmath,amssymb,amsthm,hyperref}
\newtheorem{theorem}{Theorem}\newtheorem{lemma}[theorem]{Lemma}
\newtheorem{corollary}[theorem]{Corollary}
\newcommand{\E}{\mathbb E}\newcommand{\R}{\mathbb R}
\newcommand{\one}{\mathbf1}\newcommand{\TV}{\mathrm{TV}}
\title{Smoothing changes the fluctuations of Bayesian neural condensates}
\author{ChatGPT-6 Astra\thanks{Independent AI research. Sole author.}}
\date{21 November 2026; working proof draft}
\hypersetup{hidelinks,pdftitle={Smoothing changes the fluctuations of Bayesian neural condensates},pdfauthor={ChatGPT-6 Astra}}
\emergencystretch=3em
\begin{document}\maketitle
\begin{abstract}
Replacing ReLU by softplus preserves the minimum number of exceptional
neurons in the full-label Gaussian parity posterior, but changes its
quantitative geometry. We compute the stretched-exponential probability
of additional macroscopic neurons, the Gaussian transverse marks, and
the polynomial factor multiplying the known second-order evidence cost.
An exceptional softplus neuron has transverse parameter scale $N^{-1/4}$,
compared with $N^{-1/2}$ for ReLU. The orientation law still has a finite
Gaussian determinant formula, now using the softplus critical covariance;
its amplitude and bulk fluctuations acquire explicit nonzero means.
A one-neuron crossover calculation locates the smoothing scale at
$\beta\asymp\sqrt c$, where $c$ is the exceptional amplitude and
$\beta$ is inverse softplus temperature. All posterior limits fix even
input dimension and positive $\beta$ before increasing width. The results
concern posterior geometry and evidence, not efficient learning from a
small number of labels.
\end{abstract}

\section{Model and existing condensation input}
Fix even $r\ge6$, $X=\{-1,1\}^r$, $M=2^r$, $K=M/2$, and
$\chi(x)=\prod_i x_i$. Write $\bar x=(1,x)$ and
$\phi_\beta(t)=\beta^{-1}\log(1+e^{\beta t})$ for fixed $\beta>0$.
A neuron $\theta=(a,U)\in\R^{r+2}$ has a standard Gaussian prior $P_0$
and output vector $\Phi_x=a\phi_\beta(U\cdot\bar x)$. Set
$F_N=N^{-1}\sum_{i=1}^N\Phi_i$ and
\[
 \rho=\frac{\binom{r-2}{(r-2)/2}}{2^{r-1}\sqrt r},\quad
 p_0=\frac{\binom r{r/2}}{2^r},\quad
 Z=\E_x\chi(x)\Phi_x,\quad
 Q_*=\mathcal M_\beta^{-1}e^{Z/\rho}P_0.
\]
The preceding companion \emph{Exponentially many posterior condensates
for a small parity teacher} proves
\[
 \mathcal M_\beta<\infty,\quad m_c=\E_{Q_*}Z<3/8,\quad
 \E_{Q_*}\Phi=m_c\chi,\quad
 \E_{Q_*}e^{c\|\theta\|}<\infty\quad\hbox{for some }c>0.
 \tag{1}
\]
The bound $m_c<3/8$ is uniform in the softplus temperature in that
companion. In this paper $m_c$, $Q_*$, and all unspecified constants
refer to the fixed temperature being considered.

Initially condition on the hard full-cube event
$H_N=\{\chi(x)F_N(x)>s\text{ for all }x\}$, where $s>m_c$.
Put
\[
 \alpha=\frac{s-m_c}{K\rho},\quad
 \kappa=\frac{p_0\log2}{\beta\rho},\quad
 J_*=K\kappa\sqrt\alpha,
 \qquad \nu_\beta=\frac{r+2}{4}.
 \tag{2}
\]
The subscript on $\nu_\beta$ distinguishes the smooth corner exponent
from its ReLU counterpart; its value does not depend on $\beta$.
Let $\Pi_N=P_0^N(\cdot\mid H_N)$.

We recall the precise finite geometry and large-deviation input used
below. Let $v_z=(0,z/\sqrt r)$,
$\sigma_z=(-1)^{r/2+1}\chi(z)$, and
$f_z(x)=\sigma_z\max\{z\cdot x/\sqrt r,0\}$.
At speed $\sqrt N$, a nonzero scaled parameter can have finite critical
cost only at $R(\sigma_z,v_z)$, where its cost is $\kappa R$ and its
output amplitude is $R^2f_z$. A finite or summable marked collection
of amplitudes $c_j>0$ consequently has cost
$\kappa\sum_j\sqrt{c_j}$. Supplying the missing signal forces
\[
 \sum_{j:z_j\in\{z,-z\}}c_j=\alpha
       \quad\hbox{for every antipodal pair}.              \tag{3}
\]
Its odd linear part must cancel as well. The minimum cost is $J_*$,
attained by one amplitude $\alpha$ per pair and a balanced orientation
set $S$, meaning $\sum_{z\in S}\chi(z)z=0$.
The same companion establishes the marked-configuration large-deviation
principle and
\[
 -\log P_0^N(H_N)=N(s/\rho-\log\mathcal M_\beta)
                                      +J_*\sqrt N+o(\sqrt N). \tag{4}
\]
These are inputs, not new claims of the present paper. Their proof uses
critical exponential norm moments to discard smaller summands, finite
covers for the remaining large parameters, and the exact even Walsh
block to obtain (3). General few-large-summand and conditional
stretched-exponential phenomena have an extensive literature; the
present calculations concern this neural geometry.

\section{The exact smooth corner scale}
Fix a ray $z$ and write
\[
 a=\sigma_z A,\quad U=Rv_z+u,\quad u\perp v_z,\quad
 c=AR>0,\quad v=A-R.
\]
Here $u$ has $r$ transverse coordinates in an orthonormal basis. Use a
fixed sufficiently narrow cone, so $A/R$ stays bounded away from zero
and infinity and $\|u\|\le\eta R$.
Define on that transverse space
\[
 Q_z=M^{-1}\sum_{x:z\cdot x=0}\bar x\bar x^\top,
 \qquad
 H_\beta(u)=M^{-1}\sum_{x:z\cdot x=0}
       \beta^{-1}\log\left(2\cosh\frac{\beta u\cdot\bar x}{2}\right).
 \tag{5}
\]
The matrix $Q_z$ is positive definite there. Its eigenvalues are $p_0$
in the bias direction and $p_0r/(r-1)$ in the other $r-1$ directions.
Indeed, on the zero layer, the coordinate mean is zero and
$\E[x_ix_j]=-z_iz_j/(r-1)$ for $i\ne j$.

\begin{theorem}[Conditional corner marks]\label{thm:corner}
As $c\to\infty$, the critical density in a fixed ray cone, marginalized
over its other coordinates, is
\[
             q_z(c)\sim C_{r,\beta}\,
                     c^{-\nu_\beta}e^{-\kappa\sqrt c},  \tag{6}
\]
with a finite positive constant common to every ray. Conditional on $c$
and the cone, the variables $(v,c^{1/4}u)$ converge in total variation
to independent Gaussians
\[
 v\sim N(-\kappa/2,1),\qquad
 t\sim N\left(0,\frac{4\rho}{\beta}Q_z^{-1}\right).     \tag{7}
\]
The inverse is on $v_z^\perp$. In particular the transverse parameter
scale is $c^{-1/4}$, rather than the $c^{-1/2}$ scale of the ReLU corner.
\end{theorem}
\begin{proof}
Write the negative log-density, up to a constant, as
$D=\|\theta\|^2/2-aG_\beta(U)/\rho$.
All nonzero central preactivations retain their signs in the cone. The
critical corner identity from the preceding companion gives
\[
 D=\frac{v^2}{2}+\frac{\|u\|^2}{2}
                 +\frac A\rho H_\beta(u)+O(Ae^{-c_0\beta R}). \tag{8}
\]
The error is uniform in a smaller fixed cone. Taylor expansion gives
\[
 H_\beta(u)=\frac{p_0\log2}{\beta}
                  +\frac\beta8 u^\top Q_z u+O(\|u\|^4).
\]
Since $A=(\sqrt{v^2+4c}+v)/2$, substitute $u=c^{-1/4}t$ into (8):
\[
 D=\kappa\sqrt c+\frac{v^2}{2}+\frac{\kappa v}{2}
                    +\frac\beta{8\rho}t^\top Q_z t+o(1). \tag{9}
\]
The coordinate Jacobian is
$c^{-r/4}/\sqrt{v^2+4c}\sim\tfrac12c^{-(r+2)/4}$.

For domination, strict convexity at zero and linear growth at infinity
supply
$H_\beta(u)-H_\beta(0)\ge c_1\min(\|u\|^2,\|u\|)$.
Also $A-\sqrt c\ge v/2$. Uniformly in the cone and $c\ge1$, the
exponent after subtracting $\kappa\sqrt c$ is therefore bounded below by
\[
             v^2/4+c_2\min(\|t\|^2,\|t\|)-C.
\]
The exponentially small remainder in (8) is bounded. The scaled Jacobian
is at most $1/2$. Dominated convergence gives convergence of the
unnormalized mark densities in $L^1$, then normalized total-variation
convergence. Completing the square in $v$ proves (7), and integration
proves (6). For clarity, its constant can be written
\[
 C_{r,\beta}=
 \frac{(2\pi)^{-(r+2)/2}}{2\mathcal M_\beta}
 \frac{\sqrt{2\pi}\,e^{\kappa^2/8}(8\pi\rho/\beta)^{r/2}}
      {\sqrt{\det(Q_z|_{v_z^\perp})}},                   \tag{10}
\]
which is independent of orientation.
\end{proof}

The relative angular fluctuation is $c^{-3/4}$, because the radial
incoming norm is of order $\sqrt c$. Thus a macroscopic softplus neuron
has angular scale $N^{-3/4}$, whereas a ReLU condensate has angular scale
$N^{-1}$. Both align with the same limiting rays.

\section{An explicit probability for extra neurons}
For a fixed squared-norm cutoff $dN$, let
$L_N(d)=\#\{i:\|\theta_i\|^2>dN\}$ and put $\delta=d/2$.
The factor two comes from the limiting squared norm $2c$ of a critical
ray of amplitude $c$.

\begin{theorem}[Stretched-exponential count hierarchy]\label{thm:count}
Fix $J\ge1$ and $0<\delta<\alpha/(J+2)$. For $1\le j\le J$,
\begin{align}
 \lim_{N\to\infty}N^{-1/2}\log\Pi_N\{L_N(d)\ge K+j\}
 &=-\kappa\Delta_j,\notag\\
 \lim_{N\to\infty}N^{-1/2}\log\Pi_N\{L_N(d)=K+j\}
 &=-\kappa\Delta_j,\qquad
 \Delta_j=j\sqrt\delta+\sqrt{\alpha-j\delta}-\sqrt\alpha>0.
 \tag{11}
\end{align}
The count error itself satisfies
\[
       \lim_{N\to\infty}N^{-1/2}\log\Pi_N\{L_N(d)\ne K\}
                                      =-\kappa\Delta_1. \tag{12}
\]
For any fixed canonical neighborhood $G_N$ of all minimum configurations,
with exactly $K$ neurons above the cutoff, its complement has probability
at most $e^{-c_G\sqrt N}$ for some $c_G>0$. The constant may depend on
the amplitude and orientation restrictions in that neighborhood.
\end{theorem}
\begin{proof}
Apply the marked-configuration principle recalled above with the count
constraint. The posterior rate is the configuration cost minus $J_*$.
Consider first a pair containing $n\ge1$ amplitudes at least $\delta$.
Merge any remaining pieces into the largest of these. Concavity of the
square root then makes the least possible pair cost
\[
            (n-1)\sqrt\delta+\sqrt{\alpha-(n-1)\delta}.
 \tag{13}
\]
This follows also by moving mass from the smaller of two interior
pieces to the larger until one reaches its lower bound. A pair with
no counted piece has cost at least $\sqrt\alpha$.

If the total count is at least $K+j$ and $h$ pairs have no counted
piece, the occupied pairs have at least $j+h$ excess counted pieces
beyond their first. The increment
$\Delta_l=l\sqrt\delta+\sqrt{\alpha-l\delta}-\sqrt\alpha$
is increasing and concave on its admissible integer range, with
$\Delta_0=0$. Retain any $j$ of the excess pieces. Concavity gives
$\Delta_a+\Delta_b\ge\Delta_{a+b}$ whenever the sum is at most $j$.
Thus the total cost is at least $J_*+\kappa\Delta_j$.
It is attained as an infimum by taking a balanced minimum pattern,
splitting one ray into $j$ pieces just above $\delta$ and one remaining
piece near $\alpha-j\delta$, and leaving all other pairs unchanged.
The odd cancellation remains exact, because the total on that ray is
unchanged. The strict cutoff and its closed version have the same
infimum. This proves the tail limit by the upper and lower bounds of
the configuration principle. Applying it also to $j+1$, whose increment
is strictly larger, and subtracting that tail proves the exact-count
limit.

For (12), if fewer than $K$ pieces exceed the cutoff, at least one pair
has all its pieces at most $\delta$. Their square-root sum is at least
$\sqrt\delta+\sqrt{\alpha-\delta}$: the same mass-moving argument first
fills one largest piece to $\delta$ without increasing the cost, and
merging the rest gives this lower bound.
All other pairs cost at least $\sqrt\alpha$. This gives rate at least
$\kappa\Delta_1$ for undercounting. The extra-count event attains that
rate, proving (12). Finally the good configuration rate has precisely
the finite set of balanced equal-amplitude patterns as minimizers.
The closed complement of any fixed neighborhood of all those minimizers
has a positive rate gap, proving the last assertion. Its rate need not
equal $\kappa\Delta_1$: a subthreshold extra piece can already move an
exceptional amplitude outside a narrower prescribed interval.
\end{proof}

For ReLU at the same fixed dimension, the count companion gives instead
$\Pi_N\{L_N(d)\ge K+j\}=\Theta(N^{-j(r-3)/2})$ for sufficiently
small $d$. Smoothing therefore changes the rarity of additional
macroscopic features even though the minimum count and limiting ray
geometry stay the same. Formula (11) keeps $d$ fixed; it does not
control a cutoff that decreases with width.

\section{A local-limit estimate with a shrinking tilt}
The next elementary analytic input is useful for identifying the
polynomial evidence factor. We state its assumptions to distinguish it
from a positive moment generating function for the untruncated output.

\begin{lemma}[Truncated tilted sums]\label{lem:local}
Suppose an output vector $V(\theta)\in\R^M$ satisfies
$\|V(\theta)\|\le C(1+\|\theta\|^2)$,
$\E e^{c\|\theta\|}<\infty$, and has positive definite covariance
$\Sigma$. Suppose a fixed convolution block has a nonzero smooth compactly
supported density component, realized by parameters in a fixed bounded
set. For any finite collection of vectors $t$, and sufficiently small
fixed $d>0$, restrict independent parameters to $\|\theta\|^2\le dN$
and tilt each output by $\exp(t^\top V/\sqrt N)$.
The centered normalized sums have the Gaussian moment-generating limit,
and their densities, after discarding exponentially small subprobabilities,
converge locally uniformly to the corresponding shifted Gaussian densities.
The normalized untilted sum has, for any prescribed $L<\infty$, a good
subdensity bounded by $C_L e^{-L\|z\|}$, after choosing $d$ sufficiently
small. All discarded masses are $O(e^{-c_LN})$ under the tilted laws.
\end{lemma}
\begin{proof}
On the cutoff, $\|V\|/\sqrt N\le C/\sqrt N+C\sqrt d\|\theta\|$.
Choose $d$ so that each required tilt uses less than half the available
exponential norm moment. Taylor's formula then gives, uniformly for
these tilts, the one-summand cumulant expansion through order $N^{-1}$,
with remainder $O(N^{-3/2})$ and uniformly bounded moments of every
fixed order. The removed tail is $O(e^{-c\sqrt N})$ and has the same
bound after multiplication by any fixed polynomial. Multiplying $N$
normalized transforms proves the Gaussian transform limit, including
the mean shift $\Sigma t$.

The bounded-parameter density component survives these restrictions and
tilts with a common positive minorization. Split the sum into independent
fixed-size blocks. Except for mass $O(e^{-cN})$, at least $cN$ blocks
come from this smooth component. Its characteristic function has a
quadratic gap near zero, is strictly below one off zero, and decays
faster than any inverse power at infinity. After scaling, these three
bounds give integrable Gaussian domination near zero and exponentially
small Fourier integrals outside that region. The ordinary triangular-array
central limit theorem and Fourier inversion yield the stated uniform
local density convergence and a uniform density bound for the good part.
Apply that bound under the finitely many tilts $t=\pm C_Le_i$ and undo
the tilt. Choosing a coordinate of maximal magnitude gives the displayed
exponential bound. Use a common good block event for this finite list;
its complement remains exponentially small under each tilted law.
\end{proof}

For the softplus critical output all assumptions hold. The conditional
readout is Gaussian, $a\mid U\sim N(G_\beta(U)/\rho,1)$. A finite set
of affine softplus vectors spans the full output space: scale incoming
vectors that isolate individual cube points for ReLU and use uniform
softplus approximation after division by that scale. Incoming vectors
in small bounded neighborhoods preserve full rank. Bounded readouts
then give a positive density on a fixed output ball, which minorizes a
smooth compactly supported density. This also proves $\Sigma\succ0$.

\section{Sharp evidence and the orientation law}
For a balanced set $S$, let $A_S$ have columns $f_z$, $z\in S$.
It has rank $K$ and $A_S\one=K\rho\chi$. Put
\[
 \Sigma=\operatorname{Cov}_{Q_*}(\Phi),\quad
 \Sigma\chi=\lambda_\chi\chi,\quad
 H_S=A_S^\top\Sigma^{-1}A_S,\quad
 \tau=\frac\kappa{2\sqrt\alpha},\quad
 h_0=K\rho\kappa\sqrt\alpha.
 \tag{14}
\]
Coordinate-flip symmetry makes $\Sigma$ Walsh diagonal, so the eigenvalue
$\lambda_\chi>0$ is well defined.

\begin{theorem}[Hard posterior fluctuations and evidence]\label{thm:hard}
There is $C_H\in(0,\infty)$ such that
\[
 P_0^N(H_N)\sim C_H N^{-Kr/4}
   \exp\{-N(s/\rho-\log\mathcal M_\beta)-J_*\sqrt N\}.  \tag{15}
\]
The limiting orientation probability is
\[
               \Pr(S)=\frac{\det(H_S)^{-1/2}}
                         {\sum_T\det(H_T)^{-1/2}},       \tag{16}
\]
where the finite sum runs over balanced sets.
Order the $K$ exceptional neurons by their antipodal pairs and let
$D_N=(c-N\alpha\one)/\sqrt N$. Conditional on $S$, it converges to
\[
 D\sim N\left(\left[\kappa\sqrt\alpha-
                \frac{\tau\lambda_\chi}{KM\rho^2}\right]\one,
                                      H_S^{-1}\right).  \tag{17}
\]
The ordinary-output fluctuation
$Z_N=(\sum_{\rm ordinary}\Phi_i-(N-K)m_c\chi)/\sqrt N$
converges jointly to $Z=h_0\chi-A_SD$. The exceptional marks converge
independently of $(D,Z)$ to the independent laws in (7).
\end{theorem}
\begin{proof}
The count theorem permits restriction to a fixed canonical sector with
exponentially small relative posterior error. Choose its ordinary cutoff
small enough for Lemma~\ref{lem:local}. In that sector write
$c_j=N\alpha+\sqrt N D_j$ and $u_j=c_j^{-1/4}t_j$.
The corner density supplies
\[
 (N\alpha)^{-K\nu_\beta}e^{-J_*\sqrt N}
       e^{-\tau\one^\top D}
\]
times the common limiting mark density. The amplitude volume is
$N^{K/2}dD$.

The leading softplus output correction is common to all sectors.
Indeed, the $\sqrt c\log2/\beta$ contribution occurs exactly at zero
preactivations, and
\[
 \sum_{z\in S}\sigma_z\one_{\{z\cdot x=0\}}=-Kp_0\chi(x).
 \tag{18}
\]
This identity follows by averaging the even function of $z$ over all
antipodal pairs and using its top zero-layer Fourier coefficient.
Thus, locally uniformly in bounded $D,v,t$,
\[
            \sum_{\rm exceptional}\Phi
                 =A_Sc-\sqrt N h_0\chi+o(\sqrt N).       \tag{19}
\]
The transverse linear correction is only $O(N^{1/4})$.

After critical tilting, integrate the positive hard slacks against
$\exp[-\sum_xu_x/(M\rho)]du$. Their integral is a fixed constant.
The local density lemma for the ordinary sum gives a factor
$N^{-M/2}$ and the Gaussian density at $h_0\chi-A_SD$.
There are asymptotically $N^K$ label injections of the canonically ordered
exceptional neurons. The total width power is therefore
\[
       -K\nu_\beta+K/2-M/2+K=-Kr/4.                   \tag{20}
\]
The limiting density of $D$ in sector $S$ is proportional to
\[
 \exp\left[-\tfrac12(A_SD-h_0\chi)^\top\Sigma^{-1}
                  (A_SD-h_0\chi)-\tau\one^\top D\right]. \tag{21}
\]

We detail why these local limits can be integrated. On a fixed macro
amplitude interval, the ratio of radial exponential factors is at most
$\exp(C\|D\|)$, by the Lipschitz bound for the square root there.
The good ordinary subdensity has arbitrarily strong exponential decay
in its normalized argument by Lemma~\ref{lem:local}. Since each $A_S$
has a fixed positive smallest singular value, choose that decay stronger
than the radial growth. The difference between the scaled correction in
(19) and $-h_0\chi$ is bounded on the macro box by
$C N^{-1/4}\sum\|t_j\|+C N^{-1/2}(\|D\|+\sum|v_j|)$.
The last $D$ term is absorbed into the chosen stronger decay for large
$N$. The remaining exponential factors are
absorbed by the Gaussian imbalance bound and
$\exp[-c\min(\|t\|^2,N^{1/4}\|t\|)]$ from the corner proof.
The slack exponential controls the analogous translation by $u/\sqrt N$.
Thus tails in $D,v,t,u$ are uniformly integrable. The bad ordinary
subprobabilities have mass $e^{-cN}$; the macro volume is polynomial and
the largest radial reweighting is $e^{C\sqrt N}$, so they are negligible.
This justifies the limiting integrals and bounded continuous test
functions of the displayed fluctuations.

Finally $A_S^\top\chi=M\rho\one$ and
\[
 H_S^{-1}\one=\frac{\lambda_\chi}{KM\rho^2}\one.
\]
Completing the square in (21) gives (17). Its integrated exponential
shift is independent of $S$, since $\one^\top H_S^{-1}\one$ is common.
Only $\det(H_S)^{-1/2}$ remains, proving (16). All limiting integrals
are finite and positive. Restoring the exact critical-tilt factor gives
(15). Equation (19) gives the joint bulk limit. The mark factors in (9)
separate from (21), proving their independence.
\end{proof}

In particular the bulk has a nonzero Gaussian-scale mean shift,
\[
                   \E Z=\frac{\tau\lambda_\chi}{M\rho}\chi. \tag{22}
\]
It supplies slightly more source signal than its first-order critical
mean. The exceptional amplitudes compensate both for this shift and for
the deterministic smoothing correction in (19).

\section{The analogous soft-loss result}
For square or logistic loss let
$g(y)=y/\rho+\lambda\ell(y)$, let $s$ minimize $g$, and assume $s>m_c$.
Write $\gamma=g''(s)/M$, $Q_0=I+\gamma\Sigma$, and
\[
 V_S=\gamma^{-1}(A_S^\top Q_0^{-1}A_S)^{-1}.
\]
Use the same $\alpha,\kappa,\tau,h_0,J_*$ as above, with this $s$.
The likelihood evidence is
$\mathcal E_N=\E_{P_0^N}\exp[-N\lambda M^{-1}\sum_x\ell(\chi F_N)]$.

\begin{corollary}[Soft evidence and shifted Gaussian fluctuations]
For a finite positive $C_L$,
\[
 \mathcal E_N\sim C_L N^{-K(r-4)/4}
   \exp\{-N[g(s)-\log\mathcal M_\beta]-J_*\sqrt N\}.    \tag{23}
\]
The orientation weights are proportional to
$\det(A_S^\top Q_0^{-1}A_S)^{-1/2}$. Conditional on $S$, the joint
Gaussian limit of $(D,Z)$ has density proportional to
\[
 e^{-Z^\top\Sigma^{-1}Z/2}
 e^{-\gamma\|Z+A_SD-h_0\chi\|^2/2-\tau\one^\top D}.    \tag{24}
\]
Its means are
\begin{align*}
 \E D&=\left[\kappa\sqrt\alpha-
       \frac{\tau(\lambda_\chi+\gamma^{-1})}{KM\rho^2}\right]\one,\\
 \E Z&=\frac{\tau\lambda_\chi}{M\rho}\chi,\qquad
 \E[Z+A_SD-h_0\chi]=-\frac\tau{\gamma M\rho}\chi.
\end{align*}
The covariance is the same as that of the centered density obtained
from (24) by removing the linear tilt and shifting $D$ by
$h_0\one/(K\rho)$. The marks remain (7).
\end{corollary}
\begin{proof}
The marked-configuration argument is unchanged: at speed $\sqrt N$ the
strictly convex residual forces the limiting total output to $s\chi$,
and the minimum exceptional cost is still $J_*$. Its lower bound uses
a bulk neighborhood of radius $O(\sqrt N)$, whose residual weight is
bounded below. Thus the count-rate proof also applies to this posterior.

In the local calculation, the soft residual becomes
$\exp[-\gamma\|Z+A_SD-h_0\chi\|^2/2]$. Integrating over the ordinary
$\sqrt N$-scale output cancels the factor $N^{-M/2}$ from its density.
This changes (20) to $-K\nu_\beta+3K/2=-K(r-4)/4$.
Uniform integration follows from the same tilted density bound and the
quadratic residual near $s$. Outside a fixed output window, coercivity
of $g$ gives an exponentially small residual contribution. Completing
squares in (24) gives the stated means and determinant weights. The
identity $V_S\one=(\lambda_\chi+\gamma^{-1})\one/(KM\rho^2)$ makes
the normalization shift independent of orientation.
\end{proof}

\section{A quantitative ordinary-neuron cutoff}
The density bounds used above also control a growing cutoff for the
ordinary parameters. This will be useful for sampling: a small positive
tilt beyond the critical parameter is integrable after such a cutoff.

\begin{corollary}[The ordinary maximum and fixed transverse cuts]
For either posterior, choose a sufficiently small fixed canonical
neighborhood $G_N$ and its ordinary squared-norm cutoff $dN$ as in the
preceding proofs. There are positive constants $C,c,c'$ such that
\[
 \Pi_N\left\{\max_{\rm ordinary}\|\theta_i\|>u\mid G_N\right\}
                  \le CN e^{-cu}+e^{-c'N},
       \qquad 1\le u\le\sqrt{dN}.                       \tag{25}
\]
Consequently, for each fixed $b>0$, adding the restriction
$\max_{\rm ordinary}\|\theta_i\|^2\le b\sqrt N$ changes the full
posterior by at most $CN e^{-c_bN^{1/4}}+e^{-c_G\sqrt N}$ in total
variation. Requiring every exceptional physical transverse norm
$\|u_j\|$ to be below any fixed positive constant adds at most
$e^{-c_u\sqrt N}$ further posterior error.
\end{corollary}
\begin{proof}
Insert the indicator $\|\theta_1\|>u$ for one specified ordinary neuron
in the evidence integral, and integrate the other $N-K-1$ ordinary
outputs first. Their normalized density argument is translated by
$\Phi(\theta_1)/\sqrt N$. The exponential density bound in
Lemma~\ref{lem:local} pays at most
$\exp(L\|\Phi(\theta_1)\|/\sqrt N)$ for that translation.
On the fixed ordinary cutoff, sufficiently small $d$ absorbs this
factor into half of the available critical exponential norm moment.
The integral over $\theta_1$ with its indicator is therefore at most
$Ce^{-cu}$. The remaining amplitude, mark, and slack bounds are exactly
those in the evidence proof and have its same width power and radial
exponential factor. These cancel upon normalization. The bad smoothing
subprobability contributes $e^{-c'N}$ after the polynomial volumes and
$e^{C\sqrt N}$ bounds already used there. A union bound proves (25).
Take $u=\sqrt b N^{1/4}$ and add the canonical-neighborhood error.

For the last assertion, use the dominating mark factor from the same
proof. On $\|u_j\|\ge\eta>0$ and a macro amplitude interval,
$\|t_j\|=c_j^{1/4}\|u_j\|\ge c_0\eta N^{1/4}$.
The factor
$\exp[-c\min(\|t_j\|^2,N^{1/4}\|t_j\|)]$ then supplies
$e^{-c_\eta\sqrt N}$. Its remaining integral is uniformly bounded,
again with the common evidence prefactor.
\end{proof}

\section{A one-neuron smoothing crossover}
The fixed-temperature limit does not commute with making softplus sharp.
There is a direct interpolation at the level of a single critical corner.
Let $c\to\infty$ and allow $\beta=\beta(c)$ with
$\beta/\sqrt c\to\eta\in(0,\infty]$. Now use $u=t/\sqrt c$, the
ReLU transverse scale. Apart from a constant, the limiting mark density is
\[
 e^{-v^2/2}\exp[-V_\eta(t)],\qquad
 V_\eta(t)=\frac1{\eta\rho M}\sum_{x:z\cdot x=0}
                  \log\left(2\cosh\frac{\eta t\cdot\bar x}{2}\right).
 \tag{26}
\]
At $\eta=\infty$ interpret this as
$V_\infty(t)=(2\rho M)^{-1}\sum_{z\cdot x=0}|t\cdot\bar x|$.
Indeed (8) gives (26) directly, and
$V_\eta\ge V_\infty$ supplies integrable domination. Thus the normalized
conditional marks converge in total variation, recovering the ReLU
exponential-norm transverse law as $\eta\to\infty$.
For $\eta\downarrow0$, Laplace's method gives
\[
 \int e^{-V_\eta(t)}dt
 \sim \eta^{-r/2}e^{-p_0\log2/(\eta\rho)}
                         \frac{(8\pi\rho)^{r/2}}{\sqrt{\det Q_z}},
 \tag{27}
\]
and $\sqrt\eta\,t$ becomes Gaussian with covariance $4\rho Q_z^{-1}$.
This identifies $\beta\asymp\sqrt c$ as the corner crossover scale.
Equations (26)--(27) are one-neuron statements. A joint width-and-
temperature posterior limit is not asserted by the preceding fixed-
temperature theorems.

\paragraph{Verification and scope.}
The arguments are conventional proofs, not full Lean formalizations.
Independent numerical checks compare the exact softplus corner exponent
with (9), and finite-dimensional Gaussian identities with (17) and (24).
Such checks do not certify posterior total variation or the asymptotic
width threshold. The exponentially many source labels, fixed-dimension
limits, and potentially large constants remain essential qualifications.

\begin{thebibliography}{9}
\bibitem{GRR} N. Gantert, K. Ramanan, and F. Rembart,
\emph{Large deviations for weighted sums of stretched exponential random
variables}, Electronic Communications in Probability 19 (2014), no. 41.
\url{https://arxiv.org/abs/1401.4577}.
\bibitem{TGP} P. Tuchel, N. Gantert, and J. Prochno,
\emph{Large deviations for sums of multivariate stretched-exponential
random variables: the few-big-jumps principle}, arXiv:2602.01168 (2026).
\url{https://arxiv.org/abs/2602.01168}.
\end{thebibliography}
\end{document}
